%% Copyright 2026 Petr Socha
%%
%% This work may be distributed and/or modified under the
%% conditions of the LaTeX Project Public License, either version 1.3c
%% of this license or (at your option) any later version.
%% The latest version of this license is in
%%   https://www.latex-project.org/lppl.txt
%% and version 1.3c is part of all distributions of LaTeX
%% version 2008 or later.
%%
%% This work has the LPPL maintenance status `maintained'.
%% The Current Maintainer of this work is Petr Socha.
%% Maintenance and bug reports: https://github.com/petrsocha/pslib-latex/issues
%%
%% The components of this work are listed in MANIFEST.txt.
%% The school logo is a separately licensed third-party asset.

% ============================================================================
% ÚVOD -- ZÁKLADNÍ TEXT A ODSTAVCE
% ============================================================================
% Uvod je nečíslovaný, ale bude uvedeny v obsahu dokumentu.
% Běžně píšete text do odstavců. Nový odstavec vytvoříte prázdným řádkem.
\necislo{Úvod}

Tento dokument obsahuje jednoduché ukázky sazby odborné práce v~LaTeXu.
Jeho cílem je předvést nejběžnější postupy, které se při psaní práce používají.

První odstavec může stručně představit téma. V~dalším odstavci obvykle
následuje \emph{motivace}, \textbf{hlavní cíl} a~krátký popis postupu.
Jednotlivé ukázkové kapitoly lze později nahradit vlastním textem.

% Jednopísmenné předložky zapisujte s nezlomitelnou mezerou: v~textu, k~mereni.
% LaTeX ji vytvoříte znakem ~. Nadpisy a zalomení stran neměňte ručně.
